\chapter{The Native Apparatus Presents The Residue}

\section{Forcing the residue to show}

The corridor has been read from many stations, its faces named, its turns billed. This part does something firmer than
reading: it puts the residue under \emph{apparatus} --- forces it to present itself, not to a passing glance from a
station, but to a bench built to make it show. A face read from a station is a view; a residue forced under apparatus is
a reading the machine compels the residue to yield. This section is about the difference, and about why each gauge has
its own bench.

To force the residue to show is to subject it to a procedure designed to leave it nowhere to hide. Reading a face is
passive: the machine stands at a station and receives what the corridor presents. Forcing is active: the machine runs the
residue through an apparatus that will not return a clean answer unless the residue is really there, so that the residue
must reveal itself or the apparatus will report its absence. This is the difference between noticing a quantity and
measuring it against an instrument built to catch it. The residue has been carried, honestly, since the third chapter;
now the machine builds the benches that make the carried residue declare itself as a definite reading rather than a
carried tail.

And here the partition returns in its most concrete form: each gauge's apparatus is different, because each gauge forces
the residue with the instrument its own language builds. A physical gauge forces the residue on a physical bench --- an
apparatus of the world, hardware that reads the residue as a physical quantity. This gauge forces the residue on a
\emph{computational} bench --- an apparatus of elaboration, procedures that read the residue as a decision, a cost, a
complexity. The residue is the same; the benches are different; and what each bench reads is that gauge's native face of
the one residue. The next two sections build the two benches --- first the physical, named and handed off, then this
gauge's own, which the following chapter turns into the book's first climax.

In the two verbs forcing is a tange the apparatus performs that a station cannot. A station tanges a face passively ---
selects what the corridor happens to present from that angle. A bench tanges actively --- runs a procedure that selects
the residue out whether it wishes to present or not, and funges the result into a reading the machine can stand behind.
Forcing the residue to show is the machine's tange made insistent: not receiving a face but extracting one, and funging
the extraction into a measurement. The two benches are two ways of insisting, and the residue answers each in that
bench's own language.

\section{The physical bench}

One of the benches is not this book's, and the honest thing is to name it as another gauge's and hand it across. The
physical gauge forces the residue on a bench of the world --- an apparatus of hardware that reads the residue as force,
as orbit, as field, as the physical quantities that gauge is native to. This gauge does not build that bench and cannot
read its dials; what it can do is mark that the bench exists, that it forces the \emph{same} residue this gauge forces,
and that its readings are the physical faces of the one corridor. This section is that marking, kept deliberately brief,
because the bench belongs to the volume that can cash it.

Name the physical bench's instruments once, as that gauge's, so the reader knows the residue is forced there too. The
physical volume runs the residue through a sequence of classical and quantum benches --- a weighing of attraction, a
closed orbit, a bracketed deflection, a relativistic wave equation, a variational solve, a proximity reading, an
emission threshold --- each an apparatus that reads the carried residue as a physical quantity. These are that gauge's
apparatus, pinned to what that gauge's construction actually carries, and this book neither builds them nor grades them;
it names them to mark that the physical bench forces the same residue, and it reads none of their dials, because it has
no physical apparatus of its own. Where that gauge reads a force, this gauge reads nothing --- it stands at a station it
cannot occupy --- and it says so, rather than pretending to a reading it has no bench to take.

The discipline in handing the physical bench off is the same discipline the partition built, and it is worth stating
plainly here, on the one bench this book is tempted to reach for. It would be easy to borrow the physical bench's
readings --- to let a force or a mass into this volume's account because the residue that yields them is the residue this
volume carries. That borrowing is exactly what the partition forbids. Each gauge reads its own bench and cedes the rest;
a computational reading that leaned on a physical bench's dial would be a reading this gauge did not earn, decorrelation
broken, the cross-check spoiled. So the physical bench is named and left alone. That it forces the same residue is the
whole reason the ending's agreement means something; that this gauge cannot read it is the whole reason the agreement is
not a tuning.

In the two verbs the physical bench is a set of tanges this gauge cannot perform. That bench tanges physical faces from
the residue --- selects, with hardware, readings this gauge has no apparatus to select --- and funges them into physical
quantities. This gauge marks those tanges as real and foreign: real, because they force the same residue; foreign,
because this gauge holds no instrument to run them. It funges nothing from the physical bench into its own account. The
bench is named, the residue acknowledged as shared, and the readings handed to the gauge that owns them --- so that this
gauge can turn, undistracted, to the bench that is its own.

\section{The computational bench}

The other bench is this book's, and this section builds it. The computational bench forces the residue not with hardware
but with \emph{elaboration} --- with the machine's own procedures for deciding, checking, and costing. Where the physical
bench reads the residue off an apparatus in the world, this bench reads it off the machine's act of computing: how much
deciding the residue takes, how a checker settles it, what its description costs. This is the gauge's native apparatus,
and the next chapter turns it into the first of the book's two climaxes. Here the bench is assembled.

Lay out what the computational bench can force the residue to show. It can force the residue to declare its
\emph{decidability}: run the residue through a decision procedure and read whether the procedure halts with a verdict,
and how deep it must go to reach one. It can force the residue to declare its \emph{complexity}: measure the resources ---
the beats, the steps --- the machine spends resolving it, in the manner the theory of computation makes precise~\cite{sipser1997, cook1974}. And it can force the residue to declare its \emph{description length}: ask for the shortest program
that generates it, the incompressible core the third chapter named, and read that length as a quantity~\cite{chaitin1974}.
Decidability, complexity, description length --- these are the computational bench's dials, and each reads a face of the
residue in the language of what it costs the machine to compute rather than what it weighs in the world.

The crucial feature of this bench, the one that makes the next chapter possible, is that its apparatus is made of the
\emph{same stuff} as the thing it measures. A physical bench is hardware, distinct from the residue it reads; the
machine's elaboration is not distinct from the residue --- it is more computation, of the kind the residue itself is made
of. This means the computational bench can be turned on the machine's \emph{own} act of computing. The bench that
measures how much deciding a residue takes can measure how much deciding \emph{itself} takes; the apparatus and the
object are the same kind of thing, so the apparatus can be pointed at itself. No physical bench can do this cleanly ---
hardware measuring hardware needs a second instrument --- but a computational bench measuring computation needs no second
instrument, because it is already the thing it measures. That self-directedness is the seam to the climax.

This is where Turing's~\cite{turing1936} machine and G\"odel's~\cite{godel1931} arithmetic press hardest on the
construction, and the bench is built with their walls in view. A computational bench that can turn on itself risks the
regress and the undecidability those results made exact: a checker checking itself, a procedure asked whether it halts,
a description describing its own length. The bench is built to force the residue to show \emph{without} falling into
those traps --- to read decidability, complexity, and description length only where they are decidable, and to report a
bound, not a point, exactly where the walls forbid a point. The computational bench is the apparatus that reads the
residue as computation, built at the edge of what computation can decide about itself, and disciplined to halt at that
edge rather than pretend past it.

In the two verbs the computational bench is the machine's apparatus turned on its own kind. It tanges the residue's
computational faces --- decidability, complexity, description length --- by procedures made of the same stuff as the
residue, and funges each into a reading in the machine's own currency. Because the tanging apparatus is itself
computation, it can be turned on the machine's own computing, and that self-direction is the funge the whole next chapter
closes: the bench measuring the machine, the machine measuring the bench, one apparatus forcing itself to show. The
computational bench is built; the next chapter points it at the machine's own act, and the compiler becomes the meter.
